Rains recovered the three Pauli axes from a diagonal correlation tensor,
and Van den Nest, Dehaene, and De Moor used this mechanism in a
minimal-support LU-to-LC criterion for qubit stabilizer states
\cite[Theorem~13]{Rains1999}
\cite[Theorem~1 and Corollary~1]{VanDenNestDehaeneDeMoor2005}.
Corollary~\ref{cor:full-weyl-cover} is a criterion of the same kind at
every local dimension.  No stabilizer, CSS, or linearity hypothesis
enters it, the states are arbitrary, and the coefficients on the
covering marginals are unconstrained.  For fully entangled stabilizer
states at \(q=2\) it specializes to case~(iv) of their Corollary~1:
their condition \(A_\omega=3\) at a minimal support says exactly that
the four stabilizer elements supported inside \(\omega\) project
bijectively onto the local Pauli basis at every party of \(\omega\),
which is the full-Weyl condition of Section~\ref{sec:rigidity}.
Proposition~\ref{prop:stabilizer-ame-support} is what verifies that
same condition at every prime power.
Their Theorem~1 is a finer criterion at \(q=2\) and we do not
generalize it.  It asks only that the three Paulis occur at each qubit
somewhere in the subgroup generated by all minimal elements, so the
axes at one party may come from different supports, while the marginal
argument here needs a single support to carry the whole local Weyl
basis.  Whether some prime-power criterion has their Theorem~1 as its
\(q=2\) case, with local labels drawn from several minimal supports and
exhausting \(\F_q^2\) at each party, we leave open.
